\documentclass[tikz,border=3.14mm]{standalone}
\usepackage{tikz-3dplot}
\usepackage{pgfplots}

\begin{document}

\pgfmathsetmacro{\rvec}{1.2}
\pgfmathsetmacro{\thetavec}{30}
\pgfmathsetmacro{\phivec}{60}

\pgfmathsetmacro{\reks}{.2}

%start tikz picture, and use the tdplot_main_coords style to implement the display 
%coordinate transformation provided by 3dplot
\tdplotsetmaincoords{80}{125}

\begin{tikzpicture}[scale=5,tdplot_main_coords]

%set up some coordinates 
%-----------------------
\coordinate (O) at (0,0,0);

%determine a coordinate (P) using (r,\theta,\phi) coordinates.  This command
%also determines (Pxy), (Pxz), and (Pyz): the xy-, xz-, and yz-projections
%of the point (P).
%syntax: \tdplotsetcoord{Coordinate name without parentheses}{r}{\theta}{\phi}
\tdplotsetcoord{P}{\rvec}{\thetavec}{\phivec}
\tdplotsetcoord{k}{\rvec/2}{\thetavec+0.1}{\phivec}
\tdplotsetcoord{ek1}{1}{45-\thetavec}{\phivec}

%draw figure contents
%--------------------

%draw the main coordinate system axes
\draw[thick,->] (0,0,0) -- (1,0,0) node[anchor=north east]{$x$};
\draw[thick,->] (0,0,0) -- (0,1,0) node[anchor=north west]{$y$};
\draw[thick,->] (0,0,0) -- (0,0,1) node[anchor=south]{$z$};

%draw a vector from origin to point (P) 
\draw[-stealth,color=red] (O) -- (P);

%wavenumber and polarization vectors
\draw[thick,color=red] (k) node[anchor=north west]{$\vec{k}$};

\draw[thick,->,color=blue] (0,0,0) -- ({\reks*cos(\phivec)*cos(\thetavec)},{\reks*sin(\phivec)*cos(\thetavec)},{-\reks*sin(\thetavec)}) node[anchor=south]{};
\draw[thick,color=blue] ({\reks*cos(\phivec)*cos(\thetavec)-0.05},{\reks*sin(\phivec)*cos(\thetavec)-0.05},{-\reks*sin(\thetavec)}) node[anchor=north west]{$\vec{e}_{\vec{k},1}$};

\draw[thick,->,color=blue] (0,0,0) -- ({-\reks*sin(\phivec)},{\reks*cos(\phivec)},{0}) node[anchor=south]{};
\draw[thick,color=blue] ({-\reks*sin(\phivec)-0.1},{\reks*cos(\phivec)-0.05},{0}) node[anchor=north west]{$\vec{e}_{\vec{k},2}$};

%\draw[thick,color=red] (1,0,1) node[anchor=north west]{$\vec{u}$};

%draw projection on xy plane, and a connecting line
\draw[dashed, color=red] (O) -- (Pxy);
\draw[dashed, color=red] (P) -- (Pxy);

%draw the angle \phi, and label it
%syntax: \tdplotdrawarc[coordinate frame, draw options]{center point}{r}{angle}{label options}{label}
\tdplotdrawarc{(O)}{0.2}{0}{\phivec}{anchor=north}{$\phi$}

%set the rotated coordinate system so the x'-y' plane lies within the
%"theta plane" of the main coordinate system
%syntax: \tdplotsetthetaplanecoords{\phi}
\tdplotsetthetaplanecoords{\phivec}

%draw theta arc and label, using rotated coordinate system
\tdplotdrawarc[tdplot_rotated_coords]{(0,0,0)}{0.55}{0.}{\thetavec}{anchor=south west}{$\theta$}

%set the rotated coordinate definition within display using a translation
%coordinate and Euler angles in the "z(\alpha)y(\beta)z(\gamma)" euler rotation convention
%syntax: \tdplotsetrotatedcoords{\alpha}{\beta}{\gamma}
\tdplotsetrotatedcoords{\phivec}{\thetavec}{0}

%translate the rotated coordinate system
%syntax: \tdplotsetrotatedcoordsorigin{point}
\tdplotsetrotatedcoordsorigin{(P)}

\end{tikzpicture}
\end{document}